\documentclass{situationist-preprint}

\title{The Thickness of the Veil of Ignorance}
\msid{POL-01-b}
\msfamily{POL-01}
\msdiscipline{Political Philosophy}
\mstype{Companion manuscript}
\mskeywords{veil of ignorance; original position; Harsanyi; maximin; ambiguity; level of development}

\begin{document}
\maketitle

\begin{abstract}
How much should the parties in the original position be denied? POL-01-a argues that justice is fairness, with a difference principle modified so that ties on the worst-off level are broken by the total, and that its argument does not depend on the exact thickness of the veil. We make the second claim precise. The veil must be thick enough to exclude knowledge of the society's level of development, since that knowledge would let the parties tailor principles to circumstances that the principles must survive. Beyond that, a thinner veil, one that admits imprecise or even equal probabilities of positions, leaves the choice of principle unchanged in the cases that matter: those in which principles differ over whether every position is guaranteed an adequate level, and those in which the worst-off level is fixed.
\end{abstract}

\section{The question}

The veil of ignorance can be drawn more or less thickly. At its thinnest, in Harsanyi's construction, the parties know everything except which position is theirs and assign each position an equal probability \citep{Harsanyi1975}. At its thickest, in Rawls's, they know neither their position nor any probability of occupying it, nor the particular circumstances of their society, including its economic situation and its level of civilization \citep[\S 24]{Rawls1971}. Between the two lie veils that admit some probabilistic information but not enough to fix a single distribution.

POL-01-a relies on the thick veil (Assumption~3.1 of POL-01-a) and claims that its conclusion does not depend on the exact thickness. Two questions must be separated. Which facts must the veil exclude for the device to do its work? And how far does the choice of principle change if the veil is thinner than Rawls's? We answer the first in Section~2 and the second in Section~3.

\section{What the veil must exclude}

The parties choose principles for the basic structure of a society that persists through time. Its level of development, which we take to be the level of primary goods that its institutions can secure for the worst-off position, changes: societies grow, stagnate and contract, and the same society passes through several levels in the lifetimes of its members.

\begin{principle}[Robustness]\label{p:robust}
Principles of justice for the basic structure must be acceptable to every reasonable person subject to them at every level of development the society may reach.
\end{principle}

Principle~\ref{p:robust} follows from the aim of the device as POL-01-a describes it: to identify principles each person could accept whatever her position, and so to model the requirement of POL-02 that principles be justifiable to every reasonable person subject to them. A person's position includes the time at which she lives and the state of her society then.

\begin{proposition}\label{p:level}
If the parties knew their society's level of development, they could choose principles that violate Principle~\ref{p:robust}.
\end{proposition}

\begin{proof}
Suppose the parties know that their society is rich enough that every feasible distribution gives the worst-off position a level above the floor of Assumption~3.3 of POL-01-a. Then the third of the conditions under which maximin reasoning is rational fails: no available principle risks an inadequate life. Parties who know this may rationally choose a principle that maximizes the total without regard to the worst-off level, such as the principle of average utility. Under that principle, should the society contract, nothing protects the worst-off position from falling below the floor. Persons who live in the contracted society cannot accept the principle from their standpoint, so Principle~\ref{p:robust} is violated.
\end{proof}

The same reasoning applies in reverse. Parties who knew their society to be very poor might choose principles that subordinate the basic liberties to growth, and those principles would be unacceptable once the society could secure both. Knowledge of the level of development lets the parties tailor principles to a stage, and principles of justice must not be tailored to a stage. The veil must therefore be at least thick enough to exclude it. Rawls's further exclusion of the parties' generation \citep[\S 24]{Rawls1971} rests on the same ground.

\section{What a thinner veil changes}

Suppose now that the veil excludes the level of development but admits probabilistic information about positions: a set of probability distributions, or in the limit a single uniform one. We distinguish three kinds of comparison between distributions $x$ and $y$, with worst-off levels $m_x, m_y$ and totals $t_x, t_y$.

\begin{definition}\label{d:cases}
A comparison is a \emph{floor case} if one of $x$ and $y$ places the worst-off position below the floor and the other does not; a \emph{tie case} if $m_x = m_y$; and an \emph{upper case} otherwise.
\end{definition}

Floor cases and tie cases are \emph{the cases that matter}: the first because they decide whether every position is guaranteed an adequate life, which is what the parties care about most (Assumption~3.3 of POL-01-a); the second because they are where the modified principle differs from strict maximin (Proposition~5.5 of POL-01-a).

\begin{proposition}\label{p:tie}
In tie cases, a party behind Harsanyi's thin veil, maximizing expected primary goods with equal probabilities over the $n$ positions, ranks $x$ and $y$ exactly as the modified difference principle does.
\end{proposition}

\begin{proof}
The expected level under equal probabilities is $t/n$. With $m_x = m_y$, the modified principle (Definition~5.1 of POL-01-a) ranks $x$ above $y$ if and only if $t_x > t_y$, that is, if and only if $t_x/n > t_y/n$.
\end{proof}

\begin{proposition}\label{p:floor}
In floor cases, a party behind any veil that excludes the level of development and admits at most imprecise probabilities of positions, and whose attitude to risk lies within the reasonable range for choices determining a whole life, prefers the distribution that keeps every position above the floor.
\end{proposition}

\begin{proof}
If the probabilities are imprecise, the set of admissible distributions includes some that give the position below the floor non-negligible weight. A party who evaluates options by their worst expected value over the admissible set \citep{Gilboa1989} then gives the shortfall below the floor decisive weight whenever, as in floor cases, the loss there is a loss of an adequate life and the gain above is, by Assumption~3.3 of POL-01-a, a gain over a level already adequate. If the probabilities are precise and equal, the conclusion follows from the risk-weighted framework of \citet{Buchak2017}: for choices made once, for a whole life, and for others as well as oneself, the reasonable attitude to risk weights the worst outcomes heavily enough that the loss of adequacy is not compensated by gains above the floor.
\end{proof}

Propositions~\ref{p:tie} and~\ref{p:floor} show that the thinner veil leaves the choice unchanged in the cases that matter. It may change the verdict in upper cases, where both distributions are above the floor and the one with the higher worst-off level has the lower total. There the thick veil is needed, and its justification is the one given in Chapter~6 of POL-01-a: a principle accepted because one is likely to do well under it is not thereby justifiable to the person who does badly.

\section{Objections}

\begin{objection}
Excluding the level of development leaves the parties unable to choose principles suited to their circumstances, which is what principles are for.
\end{objection}

\begin{reply}
Principles of justice are not policies. The application of the modified principle depends on the level of development at every stage, since what maximizes the worst-off level depends on what is feasible. What is excluded is knowledge that would let the choice of the principle itself depend on the stage.
\end{reply}

\begin{objection}
If the thin veil agrees with the thick one in the cases that matter, the thick veil is idle.
\end{objection}

\begin{reply}
It is not idle in upper cases, and Proposition~\ref{p:floor} for precise probabilities depends on a claim about reasonable risk attitudes that the thick veil does not need. The point is that the conclusion of POL-01 is not hostage to the exact thickness, not that thickness never matters.
\end{reply}

\section{Conclusion}

The veil must exclude the level of development, because principles chosen with that knowledge could be tailored to a stage the society will leave. Thinner than Rawls's in other respects, it would still yield the same choice where the floor is at stake and where the worst-off level is fixed. The modified difference principle of POL-01 is therefore robust to the main dispute about the veil.

\begin{thebibliography}{9}
\bibitem[Buchak(2017)]{Buchak2017} Buchak, L. (2017). Taking risks behind the veil of ignorance. \emph{Ethics}, 127(3), 610--644.
\bibitem[Gilboa and Schmeidler(1989)]{Gilboa1989} Gilboa, I. and Schmeidler, D. (1989). Maxmin expected utility with non-unique prior. \emph{Journal of Mathematical Economics}, 18(2), 141--153.
\bibitem[Harsanyi(1975)]{Harsanyi1975} Harsanyi, J.~C. (1975). Can the maximin principle serve as a basis for morality? A critique of John Rawls's theory. \emph{American Political Science Review}, 69(2), 594--606.
\bibitem[Rawls(1971)]{Rawls1971} Rawls, J. (1971). \emph{A Theory of Justice}. Harvard University Press.
\end{thebibliography}

\end{document}
